\chapter{Непертурбативное насыщение и граница картины мира-кристалла}

Однопетлевой детерминант дал отрицательную моду чистоты, но уходил к
границе потери ранга. Теперь проверяется, является ли это артефактом
гауссова приближения или полный конечномерный интеграл по матричному мосту
сам создаёт конечную одноосную фазу.

\section{Полный матричный интеграл}

При фиксированном верном состоянии $R>0$ минимальный контроль имеет вид
\begin{equation}
 Z_{\rm br}(R)
 =\int_{M_3(\mathbb R)}dB\,
  e^{-\mathcal S_R(B)},
 \qquad
 \mathcal F_{\rm br}(R)=-\log Z_{\rm br}(R),
 \label{eq:v6-full-bridge-matrix-integral}
\end{equation}
где $dB$ --- плоская мера девяти вещественных компонент, а
$\mathcal S_R$ задано предыдущим проверкой. Для строго положительного $R$
квартичный член коэрцитивен, и интеграл конечен.

Вопрос состоит в поведении при приближении к чистой страте. Положим
\begin{equation}
 R=P=\operatorname{diag}(1,0,0),
 \qquad
 B_C=\begin{pmatrix}1&0\\0&C\end{pmatrix},
 \qquad C\in M_2(\mathbb R).
 \label{eq:v6-rank-one-flat-valley}
\end{equation}
Прямое вычисление даёт
\begin{equation}
 \mathcal S_P(B_C)=0
 \qquad\text{для любого }C.
 \label{eq:v6-flat-crystal-valley}
\end{equation}
Следовательно, на проекторной границе возникает четырёхмерная плоская
долина, отсутствовавшая в локальном гессиане верного состояния.

\section{Поперечный степенной счёт}

Пусть $C=tO$, где $O\in O(2)$ и $t\gg1$. Четыре флуктуации, связывающие
выделенную ось с поперечным блоком, получают массы порядка $t^2$; ещё одна
радиальная мода остаётся порядка единицы. Поэтому гауссов объём поперечных
мод ведёт себя как
\begin{equation}
 (\det H_\perp)^{-1/2}\sim t^{-4}.
\end{equation}
Но четырёхмерный объём матриц $C$ имеет радиальную меру $t^3dt$. Вдоль
долины остаётся
\begin{equation}
 t^3dt\,t^{-4}=\frac{dt}{t},
\end{equation}
то есть логарифмическая расходимость.

Машинный аудит подтверждает этот счёт: при
$B=\operatorname{diag}(1,t,t)$ гессиан имеет четыре точных нуля, одну
ограниченную положительную моду и четыре моды, растущие как $t^2$.

\begin{theorem}[Плоская мера не насыщает проекторный переход]
При состоянии-взвешенном действии и плоской мере $dB$ полный матричный
интеграл не определён на чистой проекторной страте: четырёхмерная долина
\eqref{eq:v6-rank-one-flat-valley} создаёт логарифмическую инфракрасную
расходимость в пространстве мостов. Поэтому однопетлевой уход не
устраняется автоматически непертурбативным интегрированием.
\end{theorem}

Это отрицательный результат для конкретной меры, но не запрет фазового
механизма. Он точно показывает, какого звена не хватает: микроскопическая
мера или дополнительная производная жёсткость должна подавить поперечный
блок $C$ до достижения границы.

\section{Возврат к ранней идее мира-кристалла}

В архивной проработке проекта уже присутствовала формула «Большой взрыв
как кристаллизация». Она сопровождалась недоказанными образами решётки
Лича, одной световой нити и геометрических массовых формул. Эти элементы
не наследуются.

После строгих Томов V--VI пережило другое ядро:
\begin{enumerate}
 \item симметричная среда имеет внутреннюю отрицательную моду;
 \item упорядоченная фаза несёт ориентационную орбиту $\mathbb{RP}^2$;
 \item независимые домены могут оставить топологические ежи;
 \item спинорно-хопфов мост читает два подъёма как $+15/-15$;
 \item дефект является несогласованностью переноса, а не добавленной
 частицей.
\end{enumerate}

Это близко геометрической теории дефектов, где дислокации читаются как
кручение, а дисклинации как кривизна, и к world-nematic модели
Клейнерта--Заанена; см. \path{arXiv:cond-mat/0407469} и
\path{arXiv:gr-qc/0307033}. Но эти работы не выводят проектные классы
$15$, вес $1/7$ или рождение времени. Проектная среда пока является
операторной, а не буквальной планковской решёткой.

\begin{proposition}[Статус ранней интуиции]
Идея космической кристаллизации ретроспективно получила строгие
кинематический и локально-динамический аналоги. Она ещё не получила
микроскопической меры, обеспечивающей конечное насыщение, поэтому не
является завершённой теорией Большого взрыва.
\end{proposition}

\section{Какой член способен остановить убегание}

Энтропия фон Неймана конечна на границе и не подавляет найденную долину.
Требуется барьер, расходящийся при потере ранга, например
\begin{equation}
 -\nu\log\det R,qquad \nu>0,
 \label{eq:v6-rank-loss-barrier}
\end{equation}
либо эквивалентная жёсткость полного матричного интеграла. Сингулярные
объёмные потенциалы такого типа известны в молекулярном выводе
Ball--Majumdar, но коэффициент и знак должны следовать из проектной меры,
а не переноситься из жидких кристаллов.

\begin{nogocriterion}[Запрет буквальной решётки]
Нельзя объявлять пространство физическим кристаллом, вводить планковскую
решётку или барьер \eqref{eq:v6-rank-loss-barrier} только ради стабилизации.
Следующая проверка обязана вывести меру из конечного родительского следа,
полярного якобиана и BV-фактора либо закрыть текущую ветвь.
\end{nogocriterion}

\begin{protocol}
\begin{enumerate}
 \item ранняя идея Большого взрыва как кристаллизации:
 \textbf{найдена в архиве};
 \item буквальная решётка Лича и световая нить:
 \textbf{не наследуются};
 \item world-crystal аналог в литературе:
 \textbf{существует};
 \item полный плоский матричный интеграл при $R>0$:
 \textbf{коэрцитивен};
 \item чистая проекторная граница:
 \textbf{имеет четырёхмерную плоскую долину};
 \item поперечное подавление:
 \textbf{$t^{-4}$};
 \item остаточная радиальная мера:
 \textbf{$dt/t$};
 \item непертурбативное насыщение плоской мерой:
 \textbf{нет};
 \item условный самозапуск предыдущей проверки:
 \textbf{сохраняется, но убегает};
 \item микроскопическая мера мира-кристалла:
 \textbf{не выведена};
 \item следующая проверка:
 \textbf{полярный/BV-якобиан и барьер потери ранга}.
\end{enumerate}
\end{protocol}

Машинный сертификат сохранён в
\path{s2t_v6_state_weighted_bridge_nonperturbative_saturation_gate_results.json}.